\begin{lemma}[Change lemma, mixed location: zero-divider nonclosed case]
Fix a finite loopless multigraph with its exact edge set, a palette of
$k$ colors, and a distinguished edge $e_0$.  Let $(T,\varphi_n)$ be a
positive-rung extended Tashkinov tree.  Let $\mathcal B$ be a hierarchy
base and let
\[
 T_n=T_{n,0}\subset T_{n,1}\subset\cdots\subset
 T_{n,q}\subset T_{n,q+1}=T
\]
be a good hierarchy under $\varphi_n$, where $q$ is its last divider
index and
\[
 A=T_{n,q}^{*}=\noderef{HierarchyStarredStage}(H,q).
\]
Assume that $T$ has maximum property, that $V(T)$ is not elementary,
that $\varphi_n$ is proper and leaves exactly $e_0$ uncolored, and that
the rung number is positive.  These data form a
\noderef{HierarchyCounterexample} $C$.

Assume the six lower-rung assertions of the technical theorem, uniformly
for every positive-rung maximum ETT on the same graph, palette, and
distinguished edge, with arbitrary certified input and generating
colorings and with rung number smaller than that of $T$: lower-rung
elementarity, terminal interchangeability, the two PE stability and
stored-pair assertions, the terminal PE component assertion, the
defective-endpoint order assertion, and the mod-coloring assertion.
Assume also the same-rung induction hypothesis that every maximum ETT
with a good hierarchy having fewer than $q$ dividers is elementary.
Finally assume the large-color estimate
\[
 2n+11\leq |\overline\varphi_n(T_n)|.
\]

Put $\ell=|T|-|A|$, the number of exact tail additions above $A$, and let
\[
 T(y_{\ell-1})=T\big|_{\min\{|A|+(\ell-1),|T|\}}
\]
be the exact ordered prefix supplied by \noderef{TreeSequencePrefix}.
Assume (6.6): $V(T(y_{\ell-1}))$ is elementary under $\varphi_n$.
Assume moreover that $C$ is lexicographically minimal among all
\noderef{HierarchyCounterexample}s on the fixed graph, palette, and
$e_0$ which have exactly $q$ dividers, allowing their certified input
colorings, generating colorings, bases, ETTs, and good hierarchies to
vary.  Explicitly, no such counterexample $C'$ satisfies
\[
 \noderef{HierarchyCounterexampleLess}
  \bigl(\noderef{HierarchyCounterexampleMeasure}(C'),
        \noderef{HierarchyCounterexampleMeasure}(C)\bigr).
\]

Let $\alpha\ne\beta$ be colors missing on $T(y_{\ell-1})$.  Write
\[
 i_c=\noderef{HierarchyFirstMissingPosition}(\varphi_n,T,c),
 \qquad v_c=T.\mathit{vertex}(i_c).
\]
If at least one of $\alpha,\beta$ is not missing on $A$, assume both
\[
 v_\alpha\prec v_\beta                                      \tag{a}
\]
and the exact edge-identity exclusion
\[
 \alpha\notin
 \left\{c:\exists t\in\operatorname{Fin}|T|,
       |A|\leq t<i_\beta,
       \ \varphi_n(T.\mathit{edge}(t))=c\right\}.             \tag{b}
\]
Thus (b) is precisely
$\alpha\notin\varphi_n\langle T(v_\beta)-A\rangle$; it is not
an endpoint-only abbreviation.


Assume in addition that the first $\alpha$-missing vertex lies in the
last starred stage and the first $\beta$-missing vertex lies outside it:
\[
        v_\alpha\in V(A),\qquad v_\beta\notin V(A).
\]
Assume that $q=0$, that $\{\alpha\}$ is not closed on $A$, and
that
\[
 \alpha\in\overline\varphi_n(T_n)
 \quad\hbox{or}\quad
 \{\alpha,\beta\}\cap D_n=\varnothing.
\]
Then
\[
 \noderef{TwoColorComponent}(G,\varphi_n,\alpha,\beta,
       \operatorname{inl}v_\alpha)
 =
 \noderef{TwoColorComponent}(G,\varphi_n,\alpha,\beta,
       \operatorname{inl}v_\beta).
\]
\end{lemma}
\begin{proof}
Put $H=C.\mathit{hierarchy}$ and use the notation of the statement.  In
particular, $q=0$ and
$A=\noderef{HierarchyStarredStage}(H,0)$.  For a vertex $x$, write
\[
 K_x=\noderef{TwoColorComponent}
       (G,\varphi_n,\alpha,\beta,\operatorname{inl}x).
\]

As in the exact-prefix argument of the paper, the nestedness and terminal
fields of \noderef{TashkinovHierarchy} make $A$ an ordered initial
segment of $T$.  The full tail length $\ell=|T|-|A|$ is positive:
otherwise \noderef{TreeSequenceInitialSegment} identifies $A=T$, and
the full-length instance of \noderef{TreeSequencePrefix} makes the
assumed elementary prefix equal to $T$, contrary to the
non-elementarity field of \noderef{HierarchyCounterexample}.  Hence
$A$ is an initial segment of $T(y_{\ell-1})$.  By
\noderef{ColoredVertexSetElementary}, $A$ and both constituents of its
hierarchy core are elementary.

The least-candidate clause of
\noderef{HierarchyFirstMissingPosition} says that $v_\alpha$ misses
$\alpha$ and $v_\beta$ misses $\beta$.  Their mixed locations make them
distinct.  Elementarity of $T(y_{\ell-1})$ implies that neither misses
the other color.  Thus \noderef{MissingColors},
\noderef{PresentColors}, and properness make $K_{v_\alpha}$ and
$K_{v_\beta}$ nontrivial path components as defined by
\noderef{TwoColorNontrivialPathComponent}.

Suppose, for contradiction, that $K_{v_\alpha}\ne K_{v_\beta}$, and set
\[
             \sigma_n=\varphi_n/K_{v_\beta}
\]
using \noderef{TwoColorComponentSwap}.  By
\noderef{TwoColorComponentSwapProper}, $\sigma_n$ is proper; the switch
preserves the value \textsf{none}, so the exact edge $e_0$ remains the
unique uncolored edge.  The component through $v_\alpha$ is untouched,
while $v_\beta$ now misses $\alpha$.

The branch alternative implies $\alpha\notin D_n$.  Under pair
avoidance this is immediate.  Under
$\alpha\in\overline\varphi_n(T_n)$ it follows from the definition
\noderef{TashkinovStoredColors}, which removes colors missing on the
terminal ladder tree.  If the terminal extension were SE or RE, its
constructor in \noderef{TashkinovSeriesStep} would make every terminal
missing color outside $D_n$ closed on $T_n$; in SE the selected stable
relation and its reverse from \noderef{sc2} identify the relevant
missing sets, and in RE the coloring is unchanged.  Since
\noderef{HierarchyBase} identifies $A=T_n$ in the non-PE case, this
would contradict the assumed nonclosedness of $\alpha$ on $A$.
Therefore the terminal extension is PE, and the base gives the exact
ordered union
\[
                            A=T_n\vee R_n.                     \tag{1}
\]

We need the genuine elementarity premise for the PE closure argument.
If the rung is greater than one, restrict the certified series and
maximum property through the preceding rung and apply lower assertion
(i).  If it is one, use the elementarity assertion of
\noderef{TechnicalTheoremBase}.  Thus the pre-step coloring is
elementary on $T_n$, and
\noderef{ParallelCommonMissingClosed} applies.  If $v_\alpha$ lay in
$V(T_n)\cap V(R_n)$, its missing color $\alpha$ would be common-missing
on the two constituents and would be closed on their exact union,
contrary to the branch hypothesis.  Consequently
\[
 v_\alpha\in V(T_n)\setminus V(R_n)
 \quad\hbox{or}\quad
 v_\alpha\in V(R_n)\setminus V(T_n).                           \tag{2}
\]

In the first case of (2), conclusion (ii) of
\noderef{interchange} makes $K_{v_\alpha}$ the unique
$(\alpha,\beta)$ path component meeting $T_n$; in the second case,
conclusion (i) makes it the unique such component meeting $R_n$.
Hence the distinct component $K_{v_\beta}$ is disjoint from the
corresponding constituent.  On the other constituent, elementarity of
the core, the uniqueness of the vertex $v_\alpha$ missing $\alpha$,
and the fact that $v_\beta$ is the first vertex missing $\beta$ imply
that every vertex of $K_{v_\beta}$ encountered there has both colors
present.  Switching the full component exchanges those two incident
colors and preserves the missing-color set pointwise.  No protected
incident edge has a controlled color there, since
$\{\alpha,\beta\}\cap D_n=\varnothing$ in that alternative; in the
other alternative the relevant conclusion of \noderef{interchange}
keeps the component away from the protected exact edges.  These two
constituent checks are precisely
\[
 \noderef{CombinedStableColoring}(T_n,R_n,D_n,\varphi_n,\sigma_n).
                                                                  \tag{3}
\]

Let $B=T(v_\beta)$ be the exact ordered prefix ending at the entry
$v_\beta$.  The mixed location and membership in
$T(y_{\ell-1})$ give, for $s=|B|-|A|$,
\[
                         1\leq s\leq\ell-1<\ell.               \tag{4}
\]
Before changing any ETT data, restrict the old terminal
\noderef{TashkinovClosure} to $B$.  This produces the old-color ETT
$\mathcal E_B$ with the identical ladder series, PE extension,
maximum property, hierarchy base, and zero-divider hierarchy.

The ordered-unused premise excludes $\alpha$ from every displayed edge
of $B-A$.  The first-missing definition excludes $\beta$ from every
ordinary augmentation edge before $v_\beta$.  The only possible
exception is the individual first connecting edge of the PE core.
Equation (3) handles that identity before reconstruction.
Apply \noderef{CombinedStableClosurePreserved}: it fixes all displayed
edge identities of $T_n$ and $R_n$ and makes the same ordered $R_n$ a
closure of the same exact support prefix.  Together with the two color
exclusions, every edge in the PE seed and the augmentation sequence to
$B$ keeps its value.  At each displayed augmentation witness, the
component switch either fixes the vertex or exchanges two present
colors, so the missing witness is unchanged.  The exact sequence
restricts by \noderef{TashkinovAugmentationSequenceRestrict} and remains
a sequence of \noderef{TashkinovAugmentation}s under $\sigma_n$.

Apply \noderef{Td} to the first stable relation in (3), using the stated
lower-rung package, to obtain \noderef{ModColoring}.  Now apply
\noderef{HierarchySplitter} to the old-color prefix ETT
$\mathcal E_B$, its maximum property, (3), the transported closure, and
the verified augmentation sequence.  It returns an actual
\noderef{CorrespondingExtendedTashkinovTree} on the same exact $B$,
a transported \noderef{MaximumProperty}, a base, and a hierarchy with
the same core and divider data.  No hierarchy goodness has yet been
inferred.

We construct the zero-divider \noderef{GoodHierarchy} field by field,
using the old Gamma pairs.  Pair size and disjointness are literal.
Avoidance follows from equality of all displayed prefix edges.  Unique
refresh is vacuous at zero dividers.  For the PE parallel-base field,
the two vertexwise clauses in (3) identify the missing-color sets on
$T_n$ and $R_n$ and hence preserve the stored index set, Gamma union,
and the required cardinality.

It remains to preserve the assertion that $A$ is not
$(\Gamma^0)^-$-closed.  Unfold the old witness using
\noderef{TreeColorSetMinusClosed}, \noderef{TreeColorSetClosed}, and
\noderef{BoundaryColorEdges}.  It consists of a color
$c\in\overline\varphi_n(A)-\Gamma^0$ and an individual boundary edge
$f$ of color $c$.  If $f$ is fixed, it remains a witness.  If it is
changed, then \noderef{TwoColorComponentSwap} gives
$c\in\{\alpha,\beta\}$ and places $f$ in $K_{v_\beta}$.  The color
$\beta$ is not missing on $A$ because its first missing vertex is
outside $A$, so $c=\alpha$ and $\alpha\notin\Gamma^0$.

Assume in this last case that $A$ became $(\Gamma^0)^-$-closed under
$\sigma_n$.  Pointwise missing-set preservation in (3) would make
$\alpha$ missing and singleton-closed on $A$ under $\sigma_n$.
Conclusion (iii) of \noderef{interchange}, applied with (3), would make
$\alpha,\beta$ interchangeable on $A$ under $\sigma_n$.  But the old
distinct full components $K_{v_\alpha}$ and $K_{v_\beta}$ remain the
same two distinct full components after a complete swap; the first
meets $A$ at $v_\alpha$, and the changed boundary edge $f$ makes the
second meet $A$.  They remain nontrivial path components, contradicting
\noderef{TreeInterchangeable}.  Thus the non-minus-closed assertion is
preserved.  There are no positive-stage minus-closed fields.  We have
therefore constructed the complete good witness.  Only now does
\noderef{SameGoodHierarchy} record equality of the stages, stored sets,
and Gamma pairs.

The two vertices of $B$ missing $\alpha$ show that this corresponding
ETT is non-elementary.  Hence the reconstructed data form a
\noderef{HierarchyCounterexample} $C'$ with the same zero divider count.
The exact measure is smaller.  Indeed, let
$r_0=\noderef{HierarchyPathNumber}(H)$ and let $r'_0$ be the value for
the truncated hierarchy.  If $s<r_0$, then $r'_0\leq s<r_0$; if
$r_0\leq s$, the path tail at $r_0$ restricts to $B$, so
$r'_0\leq r_0$.  In the equality case the second coordinate of
\noderef{HierarchyCounterexampleMeasure} drops from
$|T|-|T_{n,0}|$ to $|B|-|T_{n,0}|$ by (4).  Thus
\noderef{HierarchyCounterexampleLess} holds from $C'$ to $C$, contrary
to minimality.  The supposition of two components was false, proving
the desired equality.
\end{proof}

